\documentclass[a4paper,11pt]{article}
\usepackage{jheppub}
\usepackage{amsmath}
\usepackage{amssymb}
\usepackage{lineno}
\linenumbers

\title{\boldmath Gravity Field Dynamics and Bimetric Gravity:\\
Dual Tetrahedral Geometry as the Origin of the\\
Hassan--Rosen Interaction Parameters}

\author{Stephen Nelson}
\affiliation{Field Dynamics Research\\
\href{mailto:contact@fielddynamics.org}{contact@fielddynamics.org}}

\preprint{Zenodo: [DOI to be inserted]}

\abstract{
The Hassan--Rosen bimetric gravity framework couples two
dynamical spin-2 fields through an interaction potential
governed by five coefficients $\beta_n$, which have resisted
theoretical determination for over a decade. We show that
the coupling polynomial $f(K,s) = 1 + sK + K^2$, whose
geometric origin in the tetrad equation of general relativity
is established in a companion letter, uniquely locks all five
coefficients to $\{\beta_n\} = \{1, 0, 16, 0, 1\}$ without
free parameters. Three independent confirmations fix $\beta_2 = 16$:
the full rank of the Hessian of $e_2(S)$ at the
proportional background,
a geometric channel count at the dual tetrahedral boundary,
and the dimension of the Clifford algebra $\mathrm{Cl}_{1,3}$. Orientation parity fixes
$\beta_1 = \beta_3 = 0$; exchange symmetry and normalization
fix $\beta_0 = \beta_4 = 1$. The locked potential is
ghost-free by construction, inheriting the full
Hassan--Rosen constraint structure. In the decoupling limit,
the locked parameters kill the cubic and quintic Galileon
interactions, leaving a unique quartic coupling $c_3 = 24$.
The resulting bimetric screening equation produces a
distinct, testable prediction at the MOND transition scale:
$\mu(1) = 1/\varphi^2 \approx 0.382$, where $\varphi$ is
the Golden Ratio, compared to $\mu(1) = 1/2$ for the
single-metric limit --- a difference of
$\Delta\mu = 0.118$ at the transition scale.
}

\begin{document}
\maketitle

%--------------------------------------------------------------
\section{The Open Problem}
%--------------------------------------------------------------

Bimetric gravity, developed by Hassan and Rosen in
2012~\cite{HR2012a,HR2012b}, is a consistent, ghost-free
framework for coupling two independent spin-2 fields. In
the tetrad formulation, the theory contains two tetrad
fields $e^a_\mu$ and $\tilde{e}^a_\mu$, each generating its
own metric and propagating through its own Einstein--Hilbert
kinetic term. The interaction between the two tetrads must
be a generally covariant, derivative-free scalar density.
Covariance and the antisymmetry of the exterior product
together force this interaction to take the unique form:
\begin{equation}
  V = \sum_{n=0}^{4} \beta_n\, e_n(S),
  \label{eq:potential}
\end{equation}
where $e_n(S)$ are the elementary symmetric polynomials
of the frame transition matrix $S = e^{-1}\tilde{e}$, and
the five coefficients $\beta_n$ are real numbers. Hassan
and Rosen proved that this potential is free of the
Boulware--Deser ghost~\cite{BD1972} for \emph{any} choice
of the five $\beta_n$. The form of the interaction is
uniquely fixed by consistency. The values of the
coefficients are not.

For more than a decade, the $\beta_n$ have been treated
as continuous free parameters to be constrained by fitting
to cosmological data: supernova surveys, baryon acoustic
oscillations, CMB measurements, and gravitational wave
observations. The situation is analogous to possessing a
circuit whose stability has been rigorously proven, but
whose frequency has never been derived from the circuit's
own geometry. This paper derives it.

The key is the coupling polynomial $f(K,s) = 1 + sK + K^2$,
established in a companion letter~\cite{Nelson2026a} as
the unique polynomial compatible with gravitational coupling
in three spatial dimensions. That letter shows the
polynomial is the structural content of the tetrad equation
$g_{\mu\nu} = \eta_{ab}\,e^a_\mu\,e^b_\nu$ itself, read
by rank-level on the local frame space rather than by
component count. Here we
extract the consequence of that polynomial for the
bimetric interaction: its contact coupling at $s = 0$
uniquely locks all five $\beta_n$ without fitting.

%--------------------------------------------------------------
\section{The Dual Tetrahedral Field Chamber}
%--------------------------------------------------------------

To see how the polynomial determines the $\beta_n$, we
need to understand what the bimetric interaction looks
like in coordinate space.

A tetrad is a set of four vector fields spanning the
local tangent space. When the tetrad is oriented with
tetrahedral symmetry --- the natural frame for a dual
tetrahedral field structure --- its four legs separate
into one timelike and three spacelike, with the three
spatial legs pointing along the face normals of a
regular tetrahedron.
These four directions can be written as linear projections
of the local coordinates:
\begin{equation}
  v_1 = x+y+z,\quad v_2 = x-y-z,\quad
  v_3 = -x+y-z,\quad v_4 = -x-y+z.
\end{equation}
These are the three-dimensional spatial projections of
the tetrad's four legs; the timelike direction, invariant
under spatial parity, does not participate in the face
boundary structure. The four face normals satisfy one
linear dependence ($v_1 + v_2 + v_3 + v_4 = 0$), which
is the closure condition of the tetrahedron itself.
The primary field chamber (TetA) is the region where all
four projections satisfy $v_i \leq 1$:
\begin{equation}
  \mathrm{TetA}:\quad \max(v_1, v_2, v_3, v_4) \leq 1.
\end{equation}
The inverted dual chamber (TetB), required by the
closure condition established in~\cite{Nelson2026a} which
selects the opposite-orientation branch of the dual-tetrad
geometry, satisfies the opposite inequality:
\begin{equation}
  \mathrm{TetB}:\quad \max(-v_1,-v_2,-v_3,-v_4) \leq 1.
\end{equation}
Their intersection is a central octahedral region; their
union is a stellated octahedron. The companion
letter~\cite{Nelson2026a} defines the Field Origin as the
$s = 0$ contact coupling state of the polynomial --- the
evaluation point at which the propagation term vanishes and
both sectors contribute symmetrically. The octahedral
intersection is the coordinate-space realisation of this
coupling state: the region where both tetrads meet and
where the frame transition matrix $S$ is evaluated.

The critical observation is what happens at the face
boundaries. Each face of TetA is an equilateral triangle.
The three edge midpoints of that triangle define the
medial triangle, dividing the face into four congruent
sub-regions: three corner triangles and one central
inverted triangle. This is the standard medial subdivision
--- the unique partition of an equilateral triangle into
congruent triangles that preserves the symmetry of the
face and is generated by the penetration of the dual vertex.
The stellating vertex of TetB projects onto the centroid
of the face, which is also the centroid of the medial
triangle. Repeated across all four faces of TetA, the
medial construction produces exactly
$4 \times 4 = 16$ independent boundary channels. These
channels are the geometric substrate of the bimetric
interaction: each one is an independent way for the two
field sectors to communicate at the face-to-face boundary.

This is not a new geometric object. TetA and TetB are
implicit in every tetrad-based formulation of bimetric
gravity. What has not previously been extracted is the
consequence of their boundary channel structure for the
five $\beta_n$.

%--------------------------------------------------------------
\section{Locking the Parameters}
%--------------------------------------------------------------

\subsection{The Contact Coupling of the Polynomial}

The coupling polynomial $f(K,s) = 1 + sK + K^2$ evaluated
at $K = 4$ takes three structural values depending on
the state parameter $s$:
\begin{equation}
  f(4,+1) = 21,\qquad f(4,0) = 17,\qquad f(4,-1) = 13.
\end{equation}
The state parameter $s$ reflects which field sector is
active. The propagation term $sK$ changes sign between
TetA ($s = +1$) and TetB ($s = -1$) because propagation
has a direction, while the existence condition (1) and the
full bimetric interaction ($K^2$) are symmetric between
the two sectors and do not change.

At the symmetric Field Origin, $s = 0$, the propagation
term vanishes and the polynomial reduces to its
\emph{contact coupling}:
\begin{equation}
  f(4,0) = 1 + K^2 = 1 + 16 = 17.
\end{equation}
This has two terms. The constant $1$ is the
existence condition: the shared coupling point exists
and the interaction is defined. The term $K^2 = 16$
counts the pairwise face-coupling capacity: every face of
TetA paired with every face of TetB. The mapping from
the polynomial to the potential follows from rank: the
$K^0$ term is the rank-0 contribution (a scalar, matching
$e_0(S) = 1$), and the $K^2$ term is the rank-2
contribution (a bilinear pairing, matching the pairwise
structure of $e_2(S)$). The rank-dimension $K^2 = 16$
counts the full space of ordered frame-direction pairs;
$e_2(S)$ is the unique second-order invariant of the
transition matrix that acts on this 16-dimensional space
with full sensitivity (Confirmation~1 below establishes
this through its Hessian rank). The propagation term $sK$, being
parity-odd, has no counterpart in the symmetric potential
at $s = 0$. In the bimetric
potential~\eqref{eq:potential}, the existence condition
maps to $\beta_0 e_0(S)$ and the pairwise interaction
maps to $\beta_2 e_2(S)$. Determining the value of each
coefficient is the task of the next subsection.

\subsection{Symmetry Reductions: Four Parameters Fixed}
\label{sec:sym}

Before determining the interaction weight, we fix four
of the five $\beta_n$ by symmetry.

\paragraph{Orientation parity ($\beta_1 = \beta_3 = 0$).}
The closure condition established in~\cite{Nelson2026a}
requires the two tetrahedra to carry opposite spatial
orientations. The coupling polynomial's parity structure,
also established in~\cite{Nelson2026a}, requires the
propagation term $sK$ to vanish at the symmetric Field
Origin $s = 0$: rank-1 objects are parity-odd and
contribute nothing when the two sectors are evaluated
symmetrically. In the bimetric potential, the odd-order
elementary symmetric polynomials $e_1(S)$ and $e_3(S)$
are the algebraic counterparts of this parity-odd term:
they are the terms that change under sector exchange
$S \to S^{-1}$ (via $e_n(S^{-1}) = e_{4-n}(S)$ for
$\det(S) = 1$), and therefore distinguish which tetrad
is labeled ``primary.'' At the symmetric coupling point,
any such distinguishing term must vanish, because the
physics of the interaction cannot depend on an arbitrary
labeling. Therefore $\beta_1 = \beta_3 = 0$.

\paragraph{Exchange symmetry ($\beta_0 = \beta_4$).}
The same exchange relation $e_n(S^{-1}) = e_{4-n}(S)$
used above to identify the odd-order terms also constrains
the boundary terms: it maps $e_0 \leftrightarrow e_4$.
For the interaction Lagrangian to be invariant under this
$\mathbb{Z}_2$ sector exchange, the boundary weights must
satisfy $\beta_0 = \beta_4$.

\paragraph{Normalization ($\beta_0 = 1$).}
The coupling polynomial evaluated at the Field Origin
gives $f(4,0) = 1 + K^2 = 1 + 16$. The constant term
is $K^0 = 1$: the rank-dimension of the closure
condition, established in~\cite{Nelson2026a} as the
unique rank-0 contribution to the graded frame algebra.
This fixes $\beta_0 = K^0 = 1$. Exchange symmetry then
gives $\beta_4 = \beta_0 = 1$.

\paragraph{The sole remaining parameter.}
With $\beta_1 = \beta_3 = 0$ and $\beta_0 = \beta_4 = 1$,
the interaction potential reduces to
$V = e_0(S) + \beta_2\,e_2(S) + e_4(S)$, and $\beta_2$ is
the only undetermined coefficient. The contact coupling of
the polynomial assigns $K^2 = 16$ to the pairwise
interaction term. Three independent arguments confirm that
$\beta_2 = 16$.


\subsection{Three Independent Confirmations that $\beta_2 = 16$}

Among the five symmetric polynomials $e_0$ through $e_4$
of the frame transition matrix $S$, the second-order
polynomial $e_2(S)$ is the one that counts pairwise
combinations of frame components:
\begin{equation}
  e_2(S) = \tfrac{1}{2}\left[(\mathrm{Tr}\,S)^2
  - \mathrm{Tr}(S^2)\right]
  = \sum_{a<b}\!\left(S^a_a S^b_b - S^a_b S^b_a\right).
\end{equation}

\paragraph{Confirmation~1: Hessian rank.}
The Hessian of $e_2(S)$ at the proportional background
$S = I$ is the bilinear form:
\begin{equation}
  B(H,K) = \mathrm{Tr}(H)\,\mathrm{Tr}(K)
  - \mathrm{Tr}(HK),
\end{equation}
evaluated on perturbations $H, K$ of the frame transition
matrix. We compute this Hessian explicitly on all 16
standard basis matrices $E_{ij}$ of the space of $4\times4$
matrices, producing a $16\times16$ symmetric matrix.
Gaussian elimination confirms that this matrix has full
rank~16, with no null directions. The decomposition into
irreducible subspaces confirms the result analytically:
antisymmetric matrices ($\dim = 6$) contribute positive
eigenvalues; symmetric traceless matrices ($\dim = 9$)
contribute negative eigenvalues; multiples of the identity
($\dim = 1$) contribute a positive eigenvalue. Together:
$6 + 9 + 1 = 16$, full rank.

Because $e_2(S)$ sees every independent direction of the
frame transition $S$ with nonzero sensitivity, it is a
complete quadratic observable of the frame transition.
The Hessian has 16 independent sensitivity directions.
The companion letter~\cite{Nelson2026a} establishes unit
coupling per direction as the natural normalization through
the rank-dimension $K^2 = 16$ of the graded frame algebra.
Therefore:
\begin{equation}
  \beta_2 = \mathrm{rank}\!\left(
  \nabla^2 e_2(S)\big|_{S=I}\right) = 16.
\end{equation}

\paragraph{Confirmation~2: Geometric channel count.}
As described in Section~2, the medial subdivision of the
face-to-face boundary of the dual tetrahedral chamber
produces exactly $4 \times 4 = 16$ independent coupling
channels. The companion letter~\cite{Nelson2026a}
establishes $K^2 = 16$ as the rank-dimension of the
coupling template in the graded frame algebra: the number
of independent frame-direction pairings before any values
are assigned. Each pairing contributes one structural unit
to the coupling, fixing the natural weight per channel at
unity. The weight of $e_2(S)$ must equal this channel
count. Therefore $\beta_2 = 16$.

\paragraph{Confirmation~3: Clifford algebra dimension.}
The local tetrad $e^a_\mu$ generates the Clifford algebra
$\mathrm{Cl}_{1,3}$ at each spacetime point, whose
vector-space dimension is:
\begin{equation}
  \dim(\mathrm{Cl}_{1,3}) = \sum_{k=0}^{4}\binom{4}{k}
  = 1 + 4 + 6 + 4 + 1 = 16.
\end{equation}
The frame transition matrix $S^a_b$ lives in this
16-dimensional space, matching both the Hessian rank
and the geometric channel count. This provides an
independent algebraic confirmation that 16 is the
natural interaction dimension of the tetrad frame
transition. Therefore
$\beta_2 = \dim(\mathrm{Cl}_{1,3}) = 16$.

All three confirmations converge on the same value by
independent routes. The lock is rigid.

The complete locked parameter set is:
\begin{equation}
  \boxed{(\beta_0,\, \beta_1,\, \beta_2,\, \beta_3,\,
  \beta_4) = (1,\, 0,\, 16,\, 0,\, 1).}
  \label{eq:lock}
\end{equation}
The interaction potential is therefore uniquely determined:
\begin{equation}
  V(S) = e_0(S) + 16\,e_2(S) + e_4(S).
\end{equation}

%--------------------------------------------------------------
\section{Ghost Freedom}
%--------------------------------------------------------------

A consistent interacting spin-2 theory must eliminate the
Boulware--Deser ghost~\cite{BD1972}: an unphysical scalar
degree of freedom whose negative kinetic energy would
render the vacuum unstable. We show that the locked
potential inherits ghost freedom directly from the
structure of its construction.

The interaction Lagrangian is built exclusively from
wedge products of tetrad one-forms. Because a $4$-form
must span all four spacetime dimensions, each wedge term
contains exactly one time-like differential and three
spatial differentials. The exterior algebra enforces
$dt \wedge dt = 0$, so the lapse functions $N$ and
$\tilde{N}$ enter each individual wedge term at most
linearly:
\begin{equation}
  \frac{\partial^2\mathcal{L}_{\mathrm{int}}}{\partial N^2}
  = \frac{\partial^2\mathcal{L}_{\mathrm{int}}}
  {\partial \tilde{N}^2} = 0.
\end{equation}
This primary lapse linearity places the interaction
squarely within the Hassan--Rosen invariant class.
Hassan and Rosen proved~\cite{HR2012a,HR2012b,HR2018} that
for any potential in this class, appropriate shift
redefinitions cause the constraint algebra to close and
a secondary constraint to emerge, eliminating the ghost.
Since the locked potential $\{\beta_n\} = \{1,0,16,0,1\}$
is a specific locus within this class, ghost freedom
follows immediately. No additional proof is required.

%--------------------------------------------------------------
\section{The Galileon Bridge}
%--------------------------------------------------------------

With the interaction potential locked and proven ghost-free,
we connect it to observable phenomenology via the standard
bimetric effective field theory.

In the $\Lambda_3 = (m^2 M_\mathrm{Pl})^{1/3}$ decoupling
limit, the helicity-0 scalar sector of the locked bimetric
theory acquires Galileon-type self-interactions. The
coefficients of these Galileon interactions are fixed
functions of the $\beta_n$, following the standard
de~Rham--Gabadadze--Tolley decoupling-limit
analysis~\cite{deRham2010,deRham2011}. Here we use the
bimetric convention in which $c_2$, $c_3$, $c_4$ label the
coefficients of the Galileon interactions $\mathcal{L}_3$,
$\mathcal{L}_4$, $\mathcal{L}_5$ respectively (the cubic,
quartic, and quintic Galileons). At the proportional
background $S = I$, with the mass ratio fixed by the
symmetric locked potential ($\beta_0 = \beta_4$ implies
equal effective Planck masses for the two sectors), the
Galileon coefficients reduce to~\cite{deRham2011}:
\begin{equation}
  c_2 = \beta_1 = 0, \qquad
  c_3 = \tfrac{3}{2}\beta_2 = 24, \qquad
  c_4 = \beta_3 = 0.
\end{equation}
The factor $\tfrac{3}{2}$ in $c_3 = \tfrac{3}{2}\beta_2$
arises from the combinatorial contraction of the second
elementary symmetric polynomial $e_2(S)$ against the
helicity-0 scalar in the decoupling limit and is a
standard result~\cite{deRham2011}. We follow the Galileon
normalisation convention of~\cite{deRham2011}, in which the
quartic Galileon Lagrangian is
$\mathcal{L}_4 \propto (\partial\pi)^2[(\Box\pi)^2
- (\partial_\mu\partial_\nu\pi)^2]$.

The orientation parity condition --- which set
$\beta_1 = \beta_3 = 0$ --- has a direct and striking
consequence here: it kills the cubic and quintic Galileon
interactions entirely. Only the quartic Galileon survives,
driven by $\beta_2 = 16$ alone, with coefficient
$c_3 = 24$.

This is not a tuning. It is a structural consequence of
the dual tetrahedral geometry. The two tetrahedra couple
with opposite orientations; the Galileon sector inherits
this parity, and only the orientation-even term survives.
The 16 geometric channels at the face-to-face boundary
propagate directly into the Galileon coefficient hierarchy,
fixing $c_3 = 24$ without freedom.

In the spherically symmetric quasi-static limit, the
integrated Gauss-law equation for each Galileon
$\mathcal{L}_n$ ($n = 3, 4, 5$) yields an algebraic
equation of order $n - 2$ in the field variable
$\lambda = \pi'/r$~\cite{Vainshtein1972}. The nonlinear
mapping from the field variable $\lambda$ to the physical
screening function $\mu$ reduces this by one order:
$\mathcal{L}_3$ (quadratic in $\lambda$) gives the linear
screening equation $x(1 - \mu) = \mu$, i.e.\
$\mu = x/(x+1)$, the standard simple interpolating
function; $\mathcal{L}_4$ (cubic in $\lambda$) gives the
quadratic screening equation $x(1 - \mu)^2 = \mu$; and
$\mathcal{L}_5$ (quartic in $\lambda$) would give
$N = 3$. The anchor case is $\mathcal{L}_3 \to N = 1$:
every worker in modified gravity recognises this as the
simple interpolating function. Orientation parity, by
killing $\mathcal{L}_3$ and $\mathcal{L}_5$, uniquely
selects $N = 2$.

Two ingredients contribute to the screening equation that
follows. The Galileon decoupling limit determines \emph{which}
interaction order survives: orientation parity isolates
$\mathcal{L}_4$ and fixes the algebraic order. The
scalar-tensor action of the full
theory~\cite{Nelson2026_MNRAS} supplies the universal
acceleration scale $a_0$ that makes the screening function
$\mu(x)$ mass-independent --- depending only on the ratio
$x = |\mathbf{g}|/a_0$, the same for every system --- rather
than dependent on a source-specific Vainshtein radius.
Together they produce the $N = 2$ screening equation.

In the quasi-static, weak-field regime relevant to galactic
dynamics, this $N = 2$ screening takes the algebraic form
in terms of $x = |\mathbf{g}|/a_0$ (the dimensionless
gravitational acceleration) and $\mu(x)$ (the ratio of
baryonic to total acceleration):
\begin{equation}
  x(1 - \mu)^2 = \mu,
  \label{eq:screening}
\end{equation}
This relation is
verified to machine precision against its closed-form
solution:
\begin{equation}
  \mu(x) = \frac{\sqrt{1+4x} - 1}{\sqrt{1+4x} + 1}.
  \label{eq:mu_N2}
\end{equation}
Numerical solution of Eq.~\eqref{eq:screening} via
bisection across $x \in [10^{-3}, 10^3]$ recovers
Eq.~\eqref{eq:mu_N2} to better than $10^{-6}$ relative
error at every sampled point. The closed-form solution
can also be verified algebraically: substituting
Eq.~\eqref{eq:mu_N2} directly into
Eq.~\eqref{eq:screening} and simplifying confirms the
identity to machine precision with zero residual.

%--------------------------------------------------------------
\section{The Observational Discriminant}
%--------------------------------------------------------------

The locked bimetric theory and the standard single-metric
MOND interpolating function make identical predictions in
the deep Newtonian limit ($x \gg 1$) and the deep MOND
limit ($x \ll 1$). They diverge precisely at the
transition scale $x = 1$, where the Newtonian and MOND
accelerations are equal.

At this scale, Eq.~\eqref{eq:mu_N2} gives:
\begin{equation}
  \mu(1) = \frac{\sqrt{5} - 1}{\sqrt{5} + 1}
  = \frac{3 - \sqrt{5}}{2} = \frac{1}{\varphi^2}
  \approx 0.3820,
\end{equation}
where $\varphi = (1+\sqrt{5})/2$ is the Golden Ratio.
The standard single-metric interpolating function
$\mu(x) = x/(x+1)$ predicts $\mu(1) = 1/2 = 0.5000$.
The difference is $\Delta\mu = 0.118$ at the transition
scale --- 11.8 percentage points on a quantity that ranges
from 0 to 1. This is measurable
in the MOND transition regime of galaxy rotation curves
at radii where $|\mathbf{g}| \approx a_0$.
The two predictions agree in both asymptotic limits and
differ most at the transition scale. Existing rotation
curve analyses using the single-metric function
$\mu = x/(x+1)$ achieve good fits across a range of
galaxies; the bimetric correction is an 11.8 percentage
point effect concentrated at $x \approx 1$, within the
error bars of current data for most systems~\cite{Nelson2026_MNRAS}.
High-precision kinematic measurements at the MOND
transition radius would distinguish the two predictions.

The Golden Ratio does not appear here by design. It
emerges from the algebraic structure of the $N = 2$
screening equation, which itself follows from the
orientation parity condition that kills all Galileon
couplings except $c_3 = 24$. The geometry of the dual
tetrahedral field chamber determines the Galileon
hierarchy; the Galileon hierarchy determines the
screening equation; the screening equation produces
the Golden Ratio at the transition scale. It is a
chain of forced consequences, not a choice.

%--------------------------------------------------------------
\section{Discussion and Conclusion}
%--------------------------------------------------------------

The five free parameters of the Hassan--Rosen bimetric
potential have resisted theoretical determination since
the framework was established in 2012. This paper has
shown that they are not free. They are the contact coupling
of the coupling polynomial $f(K,s) = 1 + sK + K^2$,
evaluated at the symmetric Field Origin $s = 0$.

The parameter $\beta_2 = 16$ is established by three
independent confirmations: the full rank of the Hessian
of $e_2(S)$ at the proportional background, the geometric
channel count at the dual tetrahedral boundary
($4 \times 4 = 16$), and the dimension of the Clifford
algebra $\dim(\mathrm{Cl}_{1,3}) = 16$. Three routes,
one number. The lock is rigid.

Orientation parity kills $\beta_1$ and $\beta_3$;
exchange symmetry and normalization fix $\beta_0$ and
$\beta_4$. The complete locked set $\{1,0,16,0,1\}$
lies within the Hassan--Rosen ghost-free class by
construction and inherits its constraint structure
without further proof.

In the decoupling limit, the locked parameters kill
all Galileon interactions except the quartic term,
fixing $c_3 = 24$. The unique surviving Galileon
produces the $N = 2$ screening equation, whose
closed-form solution yields $\mu(1) = 1/\varphi^2$
at the MOND transition scale. This prediction differs
from the single-metric limit by $\Delta\mu = 0.118$ and
is testable in the transition regime of observed galaxy
rotation curves.

The bimetric parameter space has not been underdetermined
because the geometric information was absent. It has been
underdetermined because the geometry of the interaction
--- the dual tetrahedral field chamber --- was not
read for its combinatorial content. The tetrahedral frame
orientation is not a gauge choice: it is the unique frame
in which the closure structure of the $D = 3$ simplex is
manifest, as established by the simplex theorem
in~\cite{Nelson2026a}. The five coefficients
$\{\beta_n\}$ are not inputs to the theory. They are
outputs of the geometry.

%--------------------------------------------------------------
\acknowledgments
%--------------------------------------------------------------

The author acknowledges the ghost-free bimetric framework
of S.~F.~Hassan and R.~A.~Rosen, and the tetrad gravity
formalism of Cartan, Einstein, and Weyl. The Galileon
EFT analysis follows the standard decoupling-limit
treatment of de~Rham, Gabadadze, and
Tolley~\cite{deRham2010,deRham2011}.

%--------------------------------------------------------------
\begin{thebibliography}{99}
%--------------------------------------------------------------

\bibitem{HR2012a}
S.~F.~Hassan and R.~A.~Rosen,
\emph{Bimetric gravity from ghost-free massive gravity},
JHEP \textbf{02} (2012) 126.

\bibitem{HR2012b}
S.~F.~Hassan and R.~A.~Rosen,
\emph{Confirmation of the secondary constraint and absence
of ghost in massive gravity and bimetric gravity},
JHEP \textbf{04} (2012) 123.

\bibitem{HR2018}
S.~F.~Hassan and R.~A.~Rosen,
\emph{Analysis of constraints and their algebra in
bimetric theory},
JHEP \textbf{08} (2018) 182.

\bibitem{BD1972}
D.~G.~Boulware and S.~Deser,
\emph{Can gravitation have a finite range?}
Phys.\ Rev.\ D \textbf{6} (1972) 3368.

\bibitem{deRham2010}
C.~de~Rham and G.~Gabadadze,
\emph{Generalization of the Fierz--Pauli action},
Phys.\ Rev.\ D \textbf{82} (2010) 044020.

\bibitem{deRham2011}
C.~de~Rham, G.~Gabadadze and A.~J.~Tolley,
\emph{Resummation of massive gravity},
Phys.\ Rev.\ Lett.\ \textbf{106} (2011) 231101.

\bibitem{Vainshtein1972}
A.~I.~Vainshtein,
\emph{To the problem of nonvanishing gravitation mass},
Phys.\ Lett.\ B \textbf{39} (1972) 393.

\bibitem{Einstein1928}
A.~Einstein,
\emph{Riemann-Geometrie mit Aufrechterhaltung des Begriffes
des Fernparallelismus},
Sitzungsber.\ Preuss.\ Akad.\ Wiss.\ (1928) 217.

\bibitem{Nelson2026a}
S.~Nelson,
\emph{Hidden in Plain Sight Since 1928: The Gravitational
Coupling Polynomial $f(K) = K^0 + K^1 + K^2$},
companion letter (2026).

\bibitem{Nelson2026_MNRAS}
S.~Nelson,
\emph{Gravity Field Dynamics: Decoding mass from observed
rotation curves},
MNRAS, submitted (2026).

\bibitem{Nelson2026b}
S.~Nelson,
\emph{Gravity Field Dynamics: Framework Essentials},
Zenodo (2026), DOI:~[to be inserted].

\end{thebibliography}

\end{document}
